\cotefile{84}
\scanpage{19}

\[
\boxed{\delta(u)\left(\mathrel{\overset{\mathrm{déf\unsure{n}}}{=}}(\operatorname{Tr}u)^2-4N(u)\right)=-\underbrace{\delta(Q_u)}_{\substack{\text{discr. de la}\\\text{forme }Q_u}}}
\]

-- Attention au signe $-$ dans le second mb\unsure{!|!.}

On note que si $Q$ est une forme quadratique à valeurs sur le module \deleted{inv. \unsure{$\Delta$} $\Lambda$, alors}\inserted{inv. $\deleted{M}\Delta'$ \deleted{(voir not. \unsure{inv.})}}, alors \deleted{\unsure{$\varphi$}}

\[
\delta(Q)\in\Gamma\left({\deleted{M}\Delta'}^{\otimes2}\otimes L^{\otimes2}\right)
\]

est défini ainsi :

\[
\varphi_Q:E\otimes E\longrightarrow\deleted{\unsure{M}}\Delta'
\]

d'où

\[
E\longrightarrow \underline{\operatorname{Hom}}(E,M)\simeq E^\vee\otimes\deleted{M}\Delta',
\]

\deleted{\unsure{et}}d'où en prenant l'ext. \unsure{ème|âme|deux} \deleted{\unsure{le} dét} module dét $=\wedge^2$ :

\[
\underbrace{\det E}_{L}\mathbin{\unsure{-|\longrightarrow}}\underbrace{\det(E^\vee)}_{L^\vee}\otimes{\Delta'}^{\otimes2}
\]

d'où,

\[
\delta(Q)\in\Gamma\,\underbrace{L^{\vee\otimes2}{L'}^{\otimes2}}_{(\Delta'\otimes L^\vee)^{\otimes2}}\deleted{\unsure{-}}.
\]

\deleted{Mais d'autre part on a} Quand $\Delta'\simeq L$, c'est une section de $\underline{\mathcal O}$, ce qui donne un sens à la formule sur les discriminants \ldots

2.4. Soit

\[
Q:E\mathbin{\unsure{-|\longrightarrow}}\Delta'\qquad Q\in\Gamma\,\underline{\operatorname{Quad}}(E,\Delta')
\]
% archive sheet number: 18
